% !TeX root = ../navier-stokes-manuscript.tex
% ============================================================
% 2.1 Vortex Stretching
% ============================================================

The momentum equation has four pieces:
\[
  \underbrace{\partial_t u}_{\text{local acceleration}}
  +
  \underbrace{(u\cdot\nabla)u}_{\text{nonlinear transport}}
  =
  \underbrace{-\nabla p}_{\text{pressure}}
  +
  \underbrace{\nu\Delta u}_{\text{viscous diffusion}}.
\]
The central regularity problem lies in the competition between nonlinear transport and viscosity.
The hope is simple: viscosity smooths the fluid faster than nonlinear transport can sharpen it.
If that remains true at every scale, the velocity stays smooth.

\subsection{Vortex Stretching}\label{sub:ThePerfectlyStretchingVortex}

That hope first comes under pressure inside a narrow material vortex tube. The same rotating fluid is pulled apart along the tube axis and squeezed across it. The symmetric strain
\[
  S:=\frac12\bigl(\nabla u+(\nabla u)^{\mathsf T}\bigr),
\]
is the part of the velocity gradient that records this deformation. Let \(e\) point along a material segment of the tube and let \(L(t)\) be its length. When the strain extends the fluid along \(e\) at rate \(\sigma(t)>0\),
\[
  Se=\sigma(t)e,
  \qquad
  \sigma(t)>0,
  \qquad
  \frac{\dot L(t)}{L(t)}
  =
  e\cdot Se
  =
  \sigma(t).
\]
The segment lengthens at the fractional rate \(\sigma(t)\).

Because the fluid is incompressible, that gain in length is paid for across the same segment. Let \(\mathcal V\) be its fixed material volume, set \(A(t):=\mathcal V/L(t)\), and write \(r_{\mathrm{eff}}(t):=\sqrt{A(t)/\pi}\). Then
\[
  \nabla\cdot u=0,
  \qquad
  \frac{d}{dt}(A L)=0,
  \qquad
  \frac{\dot A}{A}=-\sigma(t),
  \qquad
  \frac{\dot r_{\mathrm{eff}}}{r_{\mathrm{eff}}}
  =
  -\frac{\sigma(t)}{2}.
\]
Axial stretching and radial narrowing are one motion.

The rotation carried by that same fluid is the vorticity \(\omega:=\nabla\times u\). Along the moving segment it satisfies
\[
  \frac{D\omega}{Dt}
  =
  S\omega+\nu\Delta\omega,
  \qquad
  \frac{D}{Dt}
  :=
  \partial_t+u\cdot\nabla,
\]
and when the vorticity is aligned with the stretching direction, \(\omega=|\omega|e\),
\[
  S\omega
  =
  \sigma(t)\omega,
  \qquad
  \frac12\frac{D|\omega|^2}{Dt}
  =
  \sigma(t)|\omega|^2
  +
  \nu\,\omega\cdot\Delta\omega.
\]
On any interval where the full right-hand side stays positive, the same strain that lengthens the segment and thins its core also amplifies its aligned vorticity.

If the speed stays bounded by \(U_0\) on the fixed-volume material segment \(\Omega_{\mathrm{seg}}(t)\), then
\[
  E_{\mathrm{seg}}(t)
  :=
  \frac12\int_{\Omega_{\mathrm{seg}}(t)}|u(x,t)|^2\,dx
  \leq
  \frac12U_0^2\mathcal V
  <
  \infty,
  \qquad
  \left|\frac12\bigl(\nabla u-(\nabla u)^{\mathsf T}\bigr)\right|_{\mathrm F}
  =
  \frac{|\omega|}{\sqrt2}
  \leq
  |\nabla u|_{\mathrm F}.
\]
So a bounded-speed segment can keep finite kinetic energy while the rotating part of the velocity gradient grows inside a narrowing core. The gradients rise.

\divider{\NavierStokes{} Scaling}

Once the core has narrowed, scale is already part of the problem. Fix \(t_0\) and set \(\ell:=r_{\mathrm{eff}}(t_0)\). The Navier--Stokes rescaling at that width is
\[
  u_\lambda(x,t)
  :=
  \lambda u(\lambda x,\lambda^2t),
  \qquad
  p_\lambda(x,t)
  :=
  \lambda^2p(\lambda x,\lambda^2t).
\]
For \(\lambda>1\), the corresponding vortex in \(u_\lambda\) at time \(\lambda^{-2}t_0\) has width \(\lambda^{-1}\ell\). This is a comparison solution at a finer width, not a later state of the original material tube. Differentiating gives
\[
\begin{aligned}
  \partial_tu_\lambda(x,t)
  &=
  \lambda^3(\partial_tu)(\lambda x,\lambda^2t),
  \\
  (u_\lambda\cdot\nabla)u_\lambda(x,t)
  &=
  \lambda^3\bigl((u\cdot\nabla)u\bigr)(\lambda x,\lambda^2t),
  \\
  \nabla p_\lambda(x,t)
  &=
  \lambda^3(\nabla p)(\lambda x,\lambda^2t),
  \\
  \nu\Delta u_\lambda(x,t)
  &=
  \lambda^3\nu(\Delta u)(\lambda x,\lambda^2t).
\end{aligned}
\]
All four terms acquire the same \(\lambda^3\) factor. Finer width has changed the scale of the comparison without giving viscosity any relative gain over nonlinear sharpening.

\divider{Critical Height}

Now the difficulty is no longer the equation alone. A measuring norm can grow or shrink simply because the ruler changed with scale. Let \(\Lambda:=(-\Delta)^{1/2}\). At Sobolev height \(s\), a change of variables gives
\[
  \norm{u_\lambda(t)}_{\dot H^s}^2
  =
  \lambda^{2s-1}
  \norm{u(\lambda^2t)}_{\dot H^s}^2.
\]
The factor \(\lambda^{2s-1}\) is contributed by the rescaling itself, and it disappears only at \(s=\tfrac12\). Set
\[
  H(t)
  :=
  \frac12\norm{u(t)}_{\dot H^{1/2}}^2
  =
  \frac12\norm{\Lambda^{1/2}u(t)}_2^2.
\]
At that height the ruler is silent: moving to a finer comparison neither enlarges nor diminishes the measured velocity. The levels on either side move in opposite directions,
\[
  \norm{u_\lambda(t)}_2^2
  =
  \lambda^{-1}\norm{u(\lambda^2t)}_2^2,
  \qquad
  H_\lambda(t)=H(\lambda^2t),
  \qquad
  \norm{\nabla u_\lambda(t)}_2^2
  =
  \lambda\norm{\nabla u(\lambda^2t)}_2^2.
\]
Kinetic energy falls under concentration and first-derivative energy rises, while the critical height stays fixed. Finite energy still leaves the concentration problem open, and any motion of \(H\) must come from the equation itself.

\divider{Nonlinear Production versus Viscous Dissipation}

At that fixed height the opening competition can be read without any help from the change of scale. Apply \(\Lambda^{1/2}\) to the momentum equation and pair with \(\Lambda^{1/2}u\). Write the signed nonlinear contribution as
\[
  \mathcal P_{1/2}(t)
  :=
  -\operatorname{Re}
  \pair
    {\Lambda^{1/2}u(t)}
    {\Lambda^{1/2}\bigl((u(t)\cdot\nabla)u(t)\bigr)}_{L^2}.
\]
The pressure pairing vanishes because \(\nabla\cdot u=0\), while the Laplacian contributes \(-\nu\norm{\Lambda^{3/2}u}_2^2\). What remains is the exact law
\[
  H'(t)
  +
  \nu\norm{\Lambda^{3/2}u(t)}_2^2
  =
  \mathcal P_{1/2}(t).
\]
The critical height rises only when nonlinear production exceeds viscous dissipation. Under the same rescaling,
\[
  \mathcal P_{1/2}[u_\lambda](t)
  =
  \lambda^2\mathcal P_{1/2}[u](\lambda^2t),
  \qquad
  \nu\norm{\Lambda^{3/2}u_\lambda(t)}_2^2
  =
  \lambda^2\nu\norm{\Lambda^{3/2}u(\lambda^2t)}_2^2.
\]
They still scale in exactly the same way. The competition remains level, but the common \(\lambda^2\) factor is now tied to the contracted time coordinate \(t\mapsto\lambda^{-2}t\).

\divider{The Heat Clock}

That quadratic shortening is already built into the viscous part itself. For the linear heat equation \(v_t=\nu\Delta v\),
\[
  \widehat v(\xi,t+\delta t)
  =
  e^{-\nu|\xi|^2\delta t}\widehat v(\xi,t).
\]
A structure localized to width \(r\) lives at frequencies \(|\xi|\simeq r^{-1}\), so the multiplier changes by order one when \(\nu|\xi|^2\delta t\simeq1\). The corresponding viscous comparison time is
\[
  \tau_\nu(r)
  :=
  \frac{r^2}{\nu}.
\]
This is not the lifetime of the material vortex. It is the time viscosity alone needs to change a width-\(r\) structure by order one. At the finer comparison width,
\[
  \tau_\nu(\lambda^{-1}r)
  =
  \lambda^{-2}\tau_\nu(r),
\]
so the viscous clock shortens by the same factor as the time coordinate in the full \NavierStokes{} scaling. The balance has not tilted, but every narrower test comes with less time.

\divider{The Zeno Cascade}

Repeat the narrowing through the geometric sequence
\[
  r_n:=2^{-n}r_0,
  \qquad
  \tau_n:=\frac{r_n^2}{\nu},
\]
and the associated clocks satisfy
\[
  \tau_n
  =
  \frac{r_0^2}{\nu}4^{-n},
  \qquad
  \sum_{n=0}^{\infty}\tau_n
  =
  \frac{4r_0^2}{3\nu}
  <
  \infty.
\]
So the scale--clock relation by itself does not prevent infinitely many finer comparisons from fitting inside finite time.

But that is still only a schedule of widths and clocks. To describe a candidate singularity, one actual solution would have to pass through widths
\[
  r_0>r_1>r_2>\cdots\longrightarrow0
\]
at times
\[
  t_0<t_1<t_2<\cdots\uparrow T_*<\infty,
  \qquad
  t_{n+1}-t_n\asymp\tau_n,
\]
with comparison constants independent of \(n\). The remaining time would then satisfy
\[
  T_*-t_n
  \asymp
  \sum_{k=n}^{\infty}\tau_k
  =
  \frac{4r_n^2}{3\nu}.
\]
By stage \(n\), the same velocity history would have only order \(r_n^2/\nu\) of time left before \(T_*\). Spatial concentration can fit into finite time only if that one history keeps producing its next temporal response on clocks collapsing at the same quadratic rate.